\textbf{COLUNI 2019 -- Questão 01}\par\medskip Leia o poema a seguir para responder às questões 01 e 02:

\begin{center}\textbf{Despalavra}\end{center}
\begin{verse}
Hoje eu atingi o reino das imagens, o reino da despalavra.\\
Daqui vem que todas as coisas podem ter qualidades humanas.\\
Daqui vem que todas as coisas podem ter qualidades de pássaros.\\
Daqui vem que todas as pedras podem ter qualidades de sapo.\\
Daqui vem que todos os poetas podem ter qualidades de árvore.\\
Daqui vem que os poetas podem arborizar os pássaros.\\
Daqui vem que todos os poetas podem humanizar as águas.\\
Daqui vem que os poetas devem aumentar o mundo com as suas metáforas.\\
Que os poetas podem ser pré-coisas, pré-vermes, podem ser pré-musgos.\\
Daqui vem que os poetas podem compreender o mundo sem conceitos.\\
Que os poetas podem refazer o mundo por imagens, por eflúvios, por afeto.
\end{verse}
\begin{center}(BARROS, Manoel de. Despalavra. In: \_\_\_\_\_. \textbf{Meu quintal é maior do que o mundo}. Rio de Janeiro: Alfaguara, 2015, p. 117.)\end{center}

Este poema apresenta uma característica comum à obra de Manoel de Barros que consiste em:\begin{enumerate}
\item humanizar os objetos usados no nosso dia a dia.
\item problematizar o fazer poético e a matéria de sua poesia.
\item negar a capacidade de produzir significados na poesia.
\item questionar a relação existente entre as palavras e as coisas.\end{enumerate}
